% Chapter contents: the index of the kinds of derivation-like map.
% out of the 200 rows the kind census holds (27 tokenizer fragments excluded).
% STATUS: written 2026-11-05 from ledger/encyclopaedia-kinds.tsv; corrected 2026-11-07;
% explanation of the counts and of the cut added 2026-11-09 (issue #1321).
% The mined TSV has 200 rows; 27 are tokenizer fragments and the 27 are in
% ledger/derivations-encyclopaedia-kinds-excluded.txt with their reason. The count of
% genuine kinds (count>=2) is the atlas's own, provenance `% ours: the numbered list
% of the atlas-mined kinds`. The genuine named kinds with a located citation
% (Weitzenb\"ock, Shamsuddin, Lotka--Volterra, triangular, Hasse--Schmidt, ...) are defined in their chapter.
% The number in parentheses after a name is the number of works in data/atlas.tsv whose
% title names that kind; it counts works, so it is unrelated to the number of chapters,
% sections or definitions. A title naming three kinds counts once under each of the three,
% so the numbers do not add up to a number of works. The list was cut at two works: a name
% on a single title in the whole atlas is not yet a kind the literature uses.
\chapter*{Index of the \textbf{173} kinds of derivation-like map the literature names, and the \textbf{38} more this book defines}
The titles of the works in the atlas name \textbf{173} kinds of derivation-like map, % claim-ok: a count of this book's own atlas
out of the 200 rows the kind census holds. The 27 rows that are not kinds are % claim-ok: the book's own census
fragments of the mining method, and they are listed with their reason in the
kind ledger of this book.
% ours: the numbered list below is the atlas-mined kinds
Each number after a name is the number of works whose title names that kind. A % claim-ok: the counting rule of this index
title that names three kinds counts once under each of the three, so the numbers
do not add up to a number of works.
The list was cut at two works, because a name on a single title in the whole % claim-ok: the cut-off this index uses
atlas is not yet a kind the literature uses.
A name appears twice when the literature writes it two ways, joined by a dash and % claim-ok: the spelling rule of this index
written as one word: \emph{skew derivation} (179) with \emph{skew-derivation} (10), and
\emph{lie derivation} (107) with \emph{lie-derivation} (3). Each work counts once, % claim-ok: the counting rule of this index
under the spelling its title uses.
The order is the order of the census table: by count, highest first. Kinds with % claim-ok: the ordering rule of this index
the same count are in the order the mining tool met them in the atlas, so the % claim-ok: the ordering rule of this index
order within a group of equal counts carries no meaning.

\begin{enumerate}
\item \allowbreak generalized derivation \hfill(592).
\item \allowbreak skew derivation \hfill(179).
\item \allowbreak higher derivation \hfill(145).
\item \allowbreak jordan derivation \hfill(135).
\item \allowbreak lie derivation \hfill(107).
\item \allowbreak triple derivation \hfill(98).
\item \allowbreak nilpotent derivation \hfill(85).
\item \allowbreak *-derivation \hfill(82).
\item \allowbreak local derivation \hfill(80).
\item \allowbreak 2-local derivation \hfill(75).
\item \allowbreak ($\alpha$, $\beta$)-derivation \hfill(68).
\item \allowbreak n-derivation \hfill(68).
\item \allowbreak ($\sigma$, $\tau$)-derivation \hfill(67).
\item \allowbreak inner derivation \hfill(66).
\item \allowbreak left derivation \hfill(62).
\item \allowbreak unbounded derivation \hfill(34).
\item \allowbreak symmetric bi-derivation \hfill(30).
\item \allowbreak f-derivation \hfill(30).
\item \allowbreak multiplicative derivation \hfill(23).
\item \allowbreak b-generalized derivation \hfill(23).
\item \allowbreak superderivation \hfill(22).
\item \allowbreak ternary derivation \hfill(22).
\item \allowbreak ($\theta$, $\phi$)-derivation \hfill(21).
\item \allowbreak simple derivation \hfill(20).
\item \allowbreak reverse derivation \hfill(20).
\item \allowbreak (generalized)-derivation \hfill(20).
\item \allowbreak lie-type derivation \hfill(19).
\item \allowbreak outer derivation \hfill(18).
\item \allowbreak semiderivation \hfill(18).
\item \allowbreak hasse-schmidt derivation \hfill(13).
\item \allowbreak algebraic derivation \hfill(12).
\item \allowbreak additive derivation \hfill(12).
\item \allowbreak polynomial derivation \hfill(12).
\item \allowbreak monomial derivation \hfill(12).
\item \allowbreak k-derivation \hfill(12).
\item \allowbreak generic derivation \hfill(12).
\item \allowbreak normal derivation \hfill(11).
\item \allowbreak linear derivation \hfill(11).
\item \allowbreak commuting derivation \hfill(11).
\item \allowbreak double derivation \hfill(11).
\item \allowbreak integrable derivation \hfill(10).
\item \allowbreak quadratic derivation \hfill(10).
\item \allowbreak bi-derivation \hfill(10).
\item \allowbreak (f, g)-derivation \hfill(10).
\item \allowbreak generalized semiderivation \hfill(10).
\item \allowbreak skew-derivation \hfill(10).
\item \allowbreak module derivation \hfill(9).
\item \allowbreak closed derivation \hfill(9).
\item \allowbreak order derivation \hfill(9).
\item \allowbreak positive derivation \hfill(9).
\item \allowbreak twisted derivation \hfill(9).
\item \allowbreak (m, n)-jordan derivation \hfill(9).
\item \allowbreak hom-derivation \hfill(9).
\item \allowbreak point derivation \hfill(8).
\item \allowbreak finite derivation \hfill(8).
\item \allowbreak vanishing derivation \hfill(8).
\item \allowbreak ($\alpha$, $\alpha$)-derivation \hfill(8).
\item \allowbreak permuting tri-derivation \hfill(8).
\item \allowbreak local superderivation \hfill(8).
\item \allowbreak (m, n)-derivation \hfill(8).
\item \allowbreak nonsingular derivation \hfill(7).
\item \allowbreak q-skew derivation \hfill(7).
\item \allowbreak quasi-derivation \hfill(7).
\item \allowbreak weitzenb\"ock derivation \hfill(7).
\item \allowbreak central derivation \hfill(7).
\item \allowbreak 3-derivation \hfill(7).
\item \allowbreak ($\alpha$, $\beta$, $\gamma$)-derivation \hfill(7).
\item \allowbreak p-derivation \hfill(7).
\item \allowbreak random derivation \hfill(6).
\item \allowbreak triangular derivation \hfill(6).
\item \allowbreak t-derivation \hfill(6).
\item \allowbreak jordan superderivation \hfill(6).
\item \allowbreak approximate derivation \hfill(6).
\item \allowbreak continuous derivation \hfill(6).
\item \allowbreak right derivation \hfill(6).
\item \allowbreak 2-local superderivation \hfill(6).
\item \allowbreak m-derivation \hfill(6).
\item \allowbreak coderivation \hfill(5).
\item \allowbreak homo-derivation \hfill(5).
\item \allowbreak super-derivation \hfill(5).
\item \allowbreak ($\sigma$, $\tau$)-higher derivation \hfill(5).
\item \allowbreak conformal derivation \hfill(5).
\item \allowbreak $\sigma$-derivation \hfill(5).
\item \allowbreak g-derivation \hfill(5).
\item \allowbreak anti-derivation \hfill(5).
\item \allowbreak jordan-type derivation \hfill(5).
\item \allowbreak bi-semiderivation \hfill(5).
\item \allowbreak elementary derivation \hfill(4).
\item \allowbreak semisimple derivation \hfill(4).
\item \allowbreak ($\theta$, $\theta$)-derivation \hfill(4).
\item \allowbreak generalised derivation \hfill(4).
\item \allowbreak bounded derivation \hfill(4).
\item \allowbreak ($\phi$, $\psi$)-derivation \hfill(4).
\item \allowbreak symplectic derivation \hfill(4).
\item \allowbreak partial derivation \hfill(4).
\item \allowbreak r-derivation \hfill(4).
\item \allowbreak ($\sigma$, $\tau$, $\xi$)-derivation \hfill(4).
\item \allowbreak jacobian derivation \hfill(4).
\item \allowbreak preserving derivation \hfill(4).
\item \allowbreak n-higher derivation \hfill(4).
\item \allowbreak *-lie derivation \hfill(4).
\item \allowbreak universal derivation \hfill(3).
\item \allowbreak self-adjoint derivation \hfill(3).
\item \allowbreak homogeneous derivation \hfill(3).
\item \allowbreak gamma-derivation \hfill(3).
\item \allowbreak locally-finite derivation \hfill(3).
\item \allowbreak abstract derivation \hfill(3).
\item \allowbreak higher-order derivation \hfill(3).
\item \allowbreak hyponormal derivation \hfill(3).
\item \allowbreak quasi-free derivation \hfill(3).
\item \allowbreak compact derivation \hfill(3).
\item \allowbreak braided derivation \hfill(3).
\item \allowbreak factorizable derivation \hfill(3).
\item \allowbreak shamsuddin derivation \hfill(3).
\item \allowbreak orthogonal derivation \hfill(3).
\item \allowbreak multi-derivation \hfill(3).
\item \allowbreak zero derivation \hfill(3).
\item \allowbreak symmetric derivation \hfill(3).
\item \allowbreak logarithmic derivation \hfill(3).
\item \allowbreak m-integrable derivation \hfill(3).
\item \allowbreak *-double derivation \hfill(3).
\item \allowbreak left bi-derivation \hfill(3).
\item \allowbreak combinatorial derivation \hfill(3).
\item \allowbreak cubic derivation \hfill(3).
\item \allowbreak cyclotomic derivation \hfill(3).
\item \allowbreak lie superderivation \hfill(3).
\item \allowbreak $\alpha$-generalized derivation \hfill(3).
\item \allowbreak (generalized)-skew derivation \hfill(3).
\item \allowbreak multiplicative semiderivation \hfill(3).
\item \allowbreak lie-derivation \hfill(3).
\item \allowbreak ($\alpha$, $\beta$)-reverse derivation \hfill(3).
\item \allowbreak $\theta$-derivation \hfill(3).
\item \allowbreak nil derivation \hfill(2).
\item \allowbreak stable derivation \hfill(2).
\item \allowbreak hypercentral derivation \hfill(2).
\item \allowbreak bilocal derivation \hfill(2).
\item \allowbreak weight derivation \hfill(2).
\item \allowbreak diagonalizable derivation \hfill(2).
\item \allowbreak d-derivation \hfill(2).
\item \allowbreak peirce derivation \hfill(2).
\item \allowbreak odd derivation \hfill(2).
\item \allowbreak reducible derivation \hfill(2).
\item \allowbreak closable derivation \hfill(2).
\item \allowbreak noncommuting derivation \hfill(2).
\item \allowbreak constant-free derivation \hfill(2).
\item \allowbreak integral derivation \hfill(2).
\item \allowbreak near-derivation \hfill(2).
\item \allowbreak extended derivation \hfill(2).
\item \allowbreak ($\alpha$, $\beta$)-higher derivation \hfill(2).
\item \allowbreak quartic derivation \hfill(2).
\item \allowbreak b-derivation \hfill(2).
\item \allowbreak f bi-derivation \hfill(2).
\item \allowbreak volterra derivation \hfill(2).
\item \allowbreak jordan-triple derivation \hfill(2).
\item \allowbreak generalized coderivation \hfill(2).
\item \allowbreak ternary bi-derivation \hfill(2).
\item \allowbreak ($\alpha$, $\beta$)-generalized derivation \hfill(2).
\item \allowbreak leibniz-derivation \hfill(2).
\item \allowbreak generalized superderivation \hfill(2).
\item \allowbreak graded derivation \hfill(2).
\item \allowbreak weak-local derivation \hfill(2).
\item \allowbreak lotka-volterra derivation \hfill(2).
\item \allowbreak modular derivation \hfill(2).
\item \allowbreak e-derivation \hfill(2).
\item \allowbreak generalized semi-derivation \hfill(2).
\item \allowbreak poisson derivation \hfill(2).
\item \allowbreak involutive derivation \hfill(2).
\item \allowbreak fuzzy derivation \hfill(2).
\item \allowbreak schmidt derivation \hfill(2).
\item \allowbreak n-lie derivation \hfill(2).
\item \allowbreak generalized bi-derivation \hfill(2).
\item \allowbreak bracket derivation \hfill(2).
\item \allowbreak *-reverse derivation \hfill(2).
\end{enumerate}

% ---- the kinds this book defines that no title names (2026-11-18) ----
% The owner asked where generalized Lie derivations were: defined in chapter 10
% with a located citation and a theorem, and absent from the list above, because
% that list is mined from TITLES and no title uses the name. A reader looking the
% notion up found nothing. This second list is the book's own definitions: every
% \emph{...derivation} that names a kind, checked against the list above and
% added here with the chapter that defines it. The two lists answer two
% questions, and the text says which is which.
\section*{Kinds this book defines that no title names}

The list above is mined from the titles of the works in the atlas. A kind
enters it only when at least two titles name it, so a kind this book defines,
proves theorems about and cites a source for can still be absent: no title in
the literature surveyed happens to use its name. \emph{Generalized Lie
derivation} is one, and a reader who looked it up in this index found nothing
while the definition sat in Chapter 10.

The names below are the book's own: each is defined with a located citation in
the chapter given after it, and most have a theorem of their own. They are
listed by chapter so a reader can go straight to the definition. A name that
also appears above is not repeated.

\begin{description}
  \item[Prime and semiprime rings] \emph{weak Jordan derivation}.
  \item[Jordan-type maps] \emph{Jordan left derivation}.
  \item[Basic theory] \emph{principal derivation}.
  \item[Generalized, skew and twisted derivations]
    \emph{$(\alpha,\beta)$-derivation}, \emph{$\alpha$-derivation},
    \emph{$b$-generalized skew derivation},
    \emph{generalized $(\alpha,\beta)$-derivation},
    \emph{inner $\alpha$-derivation}, \emph{$l$-generalized reverse derivation},
    \emph{$r$-generalized reverse derivation},
    \emph{generalized inner derivation}.
  \item[Higher, bi- and $n$-derivations]
    \emph{$(p,\Delta)$-variate Hasse--Schmidt derivation}, \emph{biderivation},
    \emph{permuting triderivation}.
  \item[Local maps] \emph{2-local Lie derivation}, \emph{local Lie derivation},
    \emph{weak 2-local derivation}.
  \item[Matrix and triangular algebras]
    \emph{left theta-derivation}: the derivation twisted by an
    endomorphism written $\theta$, defined in the section on triangular
    algebras.
  \item[Polynomial and Weyl algebras] \emph{Jouanolou derivation}.
  \item[Lie algebras] \emph{Lie triple derivation}.
  \item[Nonassociative algebras] \emph{$(n+1)$-ary derivation},
    \emph{left superderivation}.
  \item[Banach and operator algebras] \emph{Lie $n$-derivation},
    \emph{approximate ring derivation}, \emph{bounded point derivation}.
  \item[Cohomology] \emph{secondary derivation}.
  \item[Other systems] \emph{$(l,r)$-derivation}, \emph{$(r,l)$-derivation},
    \emph{Jordan generalized derivation}, \emph{ideal derivation},
    \emph{implicative derivation}, \emph{left-right derivation},
    \emph{right-left derivation}, \emph{type I implicative derivation},
    \emph{type II implicative derivation}.
\end{description}

\section*{Identities that determine a derivation from the structure of the ring}

A second line of work asks not what a derivation \emph{is} but what an
equation about it \emph{forces} the ring to be. The equation is a
\emph{functional identity}: it holds for all values of its variables, and the
derivation and the ring's elements are the unknowns. A single identity can then
force a map to be a derivation, an inner one, or a centralizing one, and it can
force the ring to be prime, semiprime or commutative. Chapter~\ref{ch:functional-identities}
follows this line, and the notions it uses are these.

\begin{description}
  \item[Functional identity] an equation holding for all values of its
    variables, with functions as unknowns \cite[p.~9]{BresarMartindaleChebotar2007FI}.
  \item[$d$-free set] the set of functions that never meet the derivation $d$
    in the identity, which decides whether the identity says anything about $d$
    at all \cite[p.~9]{BresarMartindaleChebotar2007FI}.
  \item[Commuting map] a map $f$ with $[f(x),x]=0$ for all $x$; on a prime ring
    it is a centralizer, and a derivation that is a commuting map is inner
    modulo the centre \cite[p.~385]{Bresar1993}.
  \item[Centralizing map] a map with $[f(x),x]\in Z(A)$ for all $x$, the weaker
    form that still forces structure.
  \item[Centralizer and multiplier] the maps an identity can force a derivation
    to be, treated in their own right.
\end{description}

The authors of this line are Bresar, Chebotar, Martindale and Beidar, whose
book is the reference the chapter follows.
